\begin{proof}[Proof of \cref{obj:lem:reflection-transfer}]
\begin{enumerate}
\item For \(p\in\{1,2\}\) and \(0<s\leq1\), \cref{obj:lem:exact-reflection-budget} gives, under the stated measurability, square-integrability, and finite restriction-norm hypotheses,
\[
\sum_{a\in\mathbb Z^p}
\left(1+\frac{\pi^2\|a\|^2}{p}\right)^s
|\widehat h(a)|^2
\leq N_{p,s}(g)^2.
\]
Thus the first assertion holds with \(c=1>0\).
% lean: CausalSmith.Experimentation.BivariateSobolevDesignCapacity.reflection_transfer

\item For \(d\geq2\), \(0<s\leq1\), and
\(m\in\mathcal H_{d,s}\), the pooled conclusion of
\cref{obj:lem:exact-reflection-budget} gives
\[
\sum_{\substack{k\in\mathbb Z^d\\
1\leq|\operatorname{supp}(k)|\leq2}}
(1+\pi^2t(k))^s|\widehat F(k)|^2
\leq1\leq3072=C_s.
\]
This proves the second assertion.
% lean: CausalSmith.Experimentation.BivariateSobolevDesignCapacity.exact_reflection_budget

\item We next establish the joint sampling law for arbitrary nonnegative
integers \(n,d\). Write
\[
\mu_{\mathrm H,d}
=\text{normalized Haar measure on }\mathbb T_2^d,
\qquad
\nu_{\mathrm H,n,d}=\mu_{\mathrm H,d}^{\otimes n}.
\]
Let the independent reflection signs and the reflected sample be
\[
\eta=(\eta_{ij})_{\substack{1\leq i\leq n\\1\leq j\leq d}}
\in\{-1,1\}^{n\times d},
\qquad
\mathbb P(\eta_{ij}=1)=\mathbb P(\eta_{ij}=-1)=\frac12,
\]
\[
\widetilde X_i=(\eta_{ij}X_{ij})_{j=1}^d
\pmod{2},
\qquad
W_i=\widetilde X_i+T.
\]
The sampling measure has total mass one because its pushforward under
\(X\) is the probability law in \cref{obj:ass:uniform-draw}.

For every nonnegative Borel function
\(\varphi:\mathbb T_2^1\to[0,\infty]\), a single reflected coordinate
satisfies
\[
\begin{aligned}
\mathbb E\varphi(\widetilde X_{ij})
&=\frac12\int_0^1\varphi(u)\,du
  +\frac12\int_0^1\varphi(2-u)\,du\\
&=\frac12\int_0^2\varphi(u)\,du
 =\int_{\mathbb T_2^1}\varphi\,d\mu_{\mathrm H,1}.
\end{aligned}
\]
Here the change of variables in the second integral combines the two
halves of the period-two circle; endpoints have Lebesgue measure zero.
Taking products over coordinates and units, using independence of the
uniform coordinates and reflection signs, gives
\[
\operatorname{Law}(\widetilde X)=\nu_{\mathrm H,n,d},
\qquad
\operatorname{Law}(\widetilde X,T)
=\nu_{\mathrm H,n,d}\otimes\mu_{\mathrm H,d}.
\]
For each \(\tau\in\mathbb T_2^d\), coordinatewise Haar translation
invariance gives
\[
\bigl(\mathbf w\mapsto
(\mathbf w_i+\tau)_{i=1}^n\bigr)_{\#}
\nu_{\mathrm H,n,d}
=\nu_{\mathrm H,n,d}.
\]
Consequently, for Borel sets
\(E_{\mathrm H}\subseteq(\mathbb T_2^d)^n\) and
\(C_{\mathrm H}\subseteq\mathbb T_2^d\),
\[
\begin{aligned}
\mathbb P(W\in E_{\mathrm H},T\in C_{\mathrm H})
&=\int_{C_{\mathrm H}}
  \int_{(\mathbb T_2^d)^n}
  \mathbf1_{E_{\mathrm H}}
  \bigl((\mathbf w_i+\tau)_{i=1}^n\bigr)
  \,d\nu_{\mathrm H,n,d}(\mathbf w)
  \,d\mu_{\mathrm H,d}(\tau)\\
&=\nu_{\mathrm H,n,d}(E_{\mathrm H})
  \mu_{\mathrm H,d}(C_{\mathrm H}).
\end{aligned}
\]
Thus
\[
\operatorname{Law}(W,T)
=\nu_{\mathrm H,n,d}\otimes\mu_{\mathrm H,d}.
\]
For \(n=0\) or \(d=0\), empty products are point masses on the
corresponding unique empty tuple, and the same product and translation
identities apply.
% lean: CausalSmith.Experimentation.BivariateSobolevDesignCapacity.reflectionLaw_map_shiftedSample

\item Now impose the hypotheses of the final assertion. Folding either
reflection of a cube coordinate recovers that coordinate:
\[
b_1(u\bmod2)=u,
\qquad
b_1(-u\bmod2)=u,
\qquad 0\leq u\leq1.
\]
For the second identity, \(u>0\) has representative \(2-u\), whose fold
is \(u\); at \(u=0\), the representative and its fold are both zero.
Since the sampled coordinates belong to the cube almost surely,
\[
b_d(\widetilde X_i)=X_i,
\qquad
F(W_i-T)=F(\widetilde X_i)=m(X_i)
\quad\text{almost surely},\qquad 1\leq i\leq n.
\]
The folding map and \(F\) are Borel measurable. Integrating the
conditional assignment law \(\pi(W)\), and then using the joint law
proved above, therefore yields
\[
\begin{aligned}
\mathbb E\left|\sum_{i=1}^n Z_i m(X_i)\right|^2
={}&
\int_{(\mathbb T_2^d)^n}\int_{\mathbb T_2^d}
\int_{\mathcal Z_n}
\left|\sum_{i=1}^n z_i F(\mathbf w_i-\tau)\right|^2
\,\pi(\mathbf w,dz)\,
d\mu_{\mathrm H,d}(\tau)\,
d\nu_{\mathrm H,n,d}(\mathbf w).
\end{aligned}
\]
Because \(F\) is real-valued, the absolute value here also equals the
complex modulus of the same sum.
% lean: CausalSmith.Experimentation.BivariateSobolevDesignCapacity.reflection_kernel_fourier_identity

\item Tonelli's theorem first interchanges the sign and shift integrals:
\[
\begin{aligned}
&\int_{\mathbb T_2^d}\int_{\mathcal Z_n}
\left|\sum_{i=1}^n z_i F(\mathbf w_i-\tau)\right|^2
\,\pi(\mathbf w,dz)\,d\mu_{\mathrm H,d}(\tau)\\
&\qquad=
\int_{\mathcal Z_n}\int_{\mathbb T_2^d}
\left|\sum_{i=1}^n z_i F(\mathbf w_i-\tau)\right|^2
\,d\mu_{\mathrm H,d}(\tau)\,\pi(\mathbf w,dz).
\end{aligned}
\]
We evaluate the inner shift integral by computing its Fourier
coefficients.

The folding normalization gives
\[
\int_{\mathbb T_2^d}|F(y)|^2\,d\mu_{\mathrm H,d}(y)
=\int_{[0,1]^d}|m(x)|^2\,dx<\infty.
\]
Indeed, fix the remaining coordinates
\(x_{\mathrm{rest}}\in[0,1]^{d-1}\) and define the first-coordinate slice by
\[
m_{\mathrm{slice},x_{\mathrm{rest}}}:[0,1]\to\mathbb R,
\qquad
m_{\mathrm{slice},x_{\mathrm{rest}}}(u_{\mathrm{slice}})
=m(u_{\mathrm{slice}},x_{\mathrm{rest}}),
\qquad u_{\mathrm{slice}}\in[0,1].
\]
The two reflected halves contribute
\[
\frac12\int_0^1
|m_{\mathrm{slice},x_{\mathrm{rest}}}(u_{\mathrm{slice}})|^2
\,du_{\mathrm{slice}}
+\frac12\int_1^2
|m_{\mathrm{slice},x_{\mathrm{rest}}}(2-u_{\mathrm{slice}})|^2
\,du_{\mathrm{slice}}
=\int_0^1
|m_{\mathrm{slice},x_{\mathrm{rest}}}(u_{\mathrm{slice}})|^2
\,du_{\mathrm{slice}}.
\]
The same change of variables applies to a slice in any coordinate.
Integrating over the remaining coordinates and iterating this identity
in each coordinate, using Tonelli's theorem, proves the displayed
\(d\)-dimensional identity.

For every
\(\mathbf w\in(\mathbb T_2^d)^n\) and \(z\in\mathcal Z_n\), define
\[
S_{\mathbf w,z}(\tau)
=\sum_{i=1}^n z_i F(\mathbf w_i-\tau),
\qquad \tau\in\mathbb T_2^d,
\]
\[
\widehat S_{\mathbf w,z}(k)
=\int_{\mathbb T_2^d}
e_{-k}(\tau)S_{\mathbf w,z}(\tau)
\,d\mu_{\mathrm H,d}(\tau),
\qquad k\in\mathbb Z^d.
\]
The map \(\tau\mapsto\mathbf w_i-\tau\) preserves Haar measure, so
\[
\int_{\mathbb T_2^d}|F(\mathbf w_i-\tau)|^2
\,d\mu_{\mathrm H,d}(\tau)
=\int_{\mathbb T_2^d}|F(y)|^2\,d\mu_{\mathrm H,d}(y)<\infty.
\]
Each summand, and hence \(S_{\mathbf w,z}\), belongs to \(L^2\).

Changing variables \(y=\mathbf w_i-\tau\) and using the character
identity
\[
e_{-k}(\mathbf w_i-y)=e_{-k}(\mathbf w_i)e_k(y)
\]
gives
\[
\begin{aligned}
\int_{\mathbb T_2^d}
e_{-k}(\tau)F(\mathbf w_i-\tau)\,d\mu_{\mathrm H,d}(\tau)
&=e_{-k}(\mathbf w_i)
  \int_{\mathbb T_2^d}e_k(y)F(y)\,d\mu_{\mathrm H,d}(y)\\
&=e_{-k}(\mathbf w_i)\widehat F(-k).
\end{aligned}
\]
Linearity over the finite sum therefore gives
\[
\widehat S_{\mathbf w,z}(k)
=\widehat F(-k)\sum_{i=1}^n z_i e_{-k}(\mathbf w_i).
\]

The normalized characters form an orthonormal basis of
\(L^2(\mu_{\mathrm H,d})\). Their orthogonality follows coordinatewise:
\[
\int_{\mathbb T_2^d}
e_k(y)\overline{e_{k'}(y)}\,d\mu_{\mathrm H,d}(y)
=
\prod_{j=1}^d
\left(\frac12\int_0^2
  \exp\bigl(\mathrm i\pi(k_j-k'_j)u\bigr)\,du\right)
=
\begin{cases}
1,&k=k',\\
0,&k\neq k',
\end{cases}
\quad k,k'\in\mathbb Z^d.
\]
Trigonometric polynomials are uniformly dense in continuous functions
on the torus, and continuous functions are dense in \(L^2\); thus
\[
\overline{\operatorname{span}\{e_k:k\in\mathbb Z^d\}}^{\,L^2}
=L^2(\mu_{\mathrm H,d}).
\]
The Hilbert-space norm identity for this orthonormal basis, followed by
the bijective reindexing \(k\mapsto-k\), now gives
\[
\begin{aligned}
\int_{\mathbb T_2^d}|S_{\mathbf w,z}(\tau)|^2
\,d\mu_{\mathrm H,d}(\tau)
&=\sum_{k\in\mathbb Z^d}|\widehat S_{\mathbf w,z}(k)|^2\\
&=\sum_{k\in\mathbb Z^d}
|\widehat F(-k)|^2
\left|\sum_{i=1}^n z_i e_{-k}(\mathbf w_i)\right|^2\\
&=\sum_{k\in\mathbb Z^d}
|\widehat F(k)|^2
\left|\sum_{i=1}^n z_i e_k(\mathbf w_i)\right|^2.
\end{aligned}
\]
All terms are nonnegative, so this identity also holds as an equality
of extended nonnegative integrals and sums.
% lean: CausalSmith.Experimentation.BivariateSobolevDesignCapacity.commonShift_parseval_lintegral

\item Substitute the shift identity into the transported expectation.
Tonelli's theorem moves the nonnegative frequency sum through both the
assignment-kernel integral and the Haar-sample integral, giving
\[
\begin{aligned}
\mathbb E\left|\sum_{i=1}^n Z_i m(X_i)\right|^2
&=
\int_{(\mathbb T_2^d)^n}\int_{\mathcal Z_n}
\sum_{k\in\mathbb Z^d}
|\widehat F(k)|^2
\left|\sum_{i=1}^n z_i e_k(\mathbf w_i)\right|^2
\,\pi(\mathbf w,dz)\,d\nu_{\mathrm H,n,d}(\mathbf w)\\
&=
\sum_{k\in\mathbb Z^d}|\widehat F(k)|^2
\int_{(\mathbb T_2^d)^n}\int_{\mathcal Z_n}
\left|\sum_{i=1}^n z_i e_k(\mathbf w_i)\right|^2
\,\pi(\mathbf w,dz)\,d\nu_{\mathrm H,n,d}(\mathbf w)\\
&=
\sum_{k\in\mathbb Z^d}|\widehat F(k)|^2
\mathbb E\left|\sum_{i=1}^n Z_i e_k(W_i)\right|^2.
\end{aligned}
\]
The last expectation uses the independent Haar sample and conditional
assignment law \(\pi(W)\), as asserted.
% lean: CausalSmith.Experimentation.BivariateSobolevDesignCapacity.commonShift_kernel_parseval
\end{enumerate}
\end{proof}
